\section{Parallel execution: from one core to a graphics processor}\label{sec:parallel}

The absolute chart times quoted so far depend on the hardware they were
measured on, and a reader with a different machine cannot transfer them. A
more useful description separates two numbers: the cost of one parameter point
on one thread, and how that cost divides among many threads. This section
measures both, on CPUs and GPUs, and shows that the chart time is proportional
to $1/n$ in the number of threads $n$ over several decades -- which is what
makes real-time charts possible, and what lets a reader estimate the chart time
on their own hardware.

\subsection{Why the sweep parallelizes}\label{sec:parallel-why}

A stability chart is an embarrassingly parallel computation: the parameter
points are independent, each needs only the user's characteristic function and
the march of Section~\ref{sec:unwrap}, and nothing is exchanged between them.
The discrete march adds three properties that a parallel machine rewards: its
state is a fixed handful of scalars (no allocation, no solver object), every
step is the same short sequence of operations (one evaluation of $\Dfun$ and
$\Dfun'$ by dual numbers, a comparison, a step update), and its cost per point
varies only moderately across a chart -- the median point of the three
benchmark systems of Table~\ref{tab:unwrap} needs $31$, $52$ and $129$
evaluations, the mean $32$, $53$ and $139$. The march was therefore written once as a
\pkg{KernelAbstractions.jl} kernel, which compiles the same source for a
multithreaded CPU and for a GPU with one thread per point; all measurements of
this section use that kernel. (The package timings of
Section~\ref{sec:performance-unwrap} use the package's own CPU implementation,
which carries the general root buffer and costs about twice as much per point:
$0.96\,\mu$s against $0.49\,\mu$s of wall time per point on the twelve
threads of the same desktop processor.)

\subsection{Measured scaling}\label{sec:parallel-scaling}

\begin{figure}[tb]
  \centering
  \maybefig{fig_parallel_scaling}{\linewidth}
  \caption{Parallel scaling of the discrete march on the fourth-order benchmark
  (Appendix~\ref{app:gallery}): time for $10^{6}$ parameter points -- so that
  $1$\,ms corresponds to $1$\,ns per point -- versus the number of threads $n$,
  on two CPUs (Float64; a six-core Intel i5-10400 desktop and the $48$-thread
  host of the Colab G4 instance) and on four NVIDIA GPUs (Google Colab), where
  exactly $n$ GPU threads march the chart (solid: Float32, dashed: Float64). On
  the GPUs the $n$ threads draw points from a shared work queue, in blocks of
  $\min(256,n)$ threads; the chart grows from $64\times64$ to
  $2048\times2048$ points with $n$, so that each thread marches at least $64$
  points up to $n=65\,536$; each value is the best of three runs. The dotted
  lines are $1/n$ through the single-thread CPU time and through the
  RTX PRO 6000 at $n=64$; the horizontal lines mark the per-point cost at which
  a full-HD or an 8K chart is computed thirty times per second. The T4 curve is
  affected by the clock control of this passively cooled card.}\label{fig:parallel}
\end{figure}

\begin{table}[tb]
  \centering
  \caption{The quantities behind Fig.~\ref{fig:parallel}: the cost of one point
  on one thread ($64\times64$ chart), the per-point cost of the saturated device
  ($2048\times2048$ chart, work-queue schedule), and their ratio, the
  \emph{effective parallelism} $P_{\mathrm{eff}}=t_1/t_{\mathrm{sat}}$.
  $^\ast$For the GPUs in Float32, $P_{\mathrm{eff}}$ divided by the number of
  CUDA cores. With one thread per point instead of the work queue the GPUs are
  $1.3$--$1.9$ times faster (Table~\ref{tab:gpu}).}\label{tab:parallel}
  \maybeinput{tables/tab_parallel.tex}
\end{table}

Figure~\ref{fig:parallel} shows the time of a chart of $10^{6}$ points as a
function of the number of threads, and Table~\ref{tab:parallel} extracts the
characteristic numbers. On the CPU the result is what the independence of the
points promises. One thread of the server needs $\ParCpuOneUs\,\mu$s per point
(Float64, i.e.\ double precision, scalar code); up to $16$ threads every doubling of the threads
almost halves the time ($\ParCpuSixteenSpeedup$ times faster at $16$), and
$48$ threads are $\ParCpuFortyEightSpeedup$ times faster -- the deficit at the
top end being consistent with simultaneous multithreading, where two hardware
threads share one core. The six-core desktop follows the same line up to its
six cores ($5.6$ times) and reaches $8.7$ times with its twelve hardware
threads. A chart that takes $3.8$\,s on one core ($10^{6}$ points of the
fourth-order benchmark) thus takes $0.1$\,s on the $48$-thread server.

On the GPU a single thread is slower than a CPU thread --
$\ParGpuOneUsLo$--$\ParGpuOneUsHi\,\mu$s per point in Float32 (single
precision), five to eight
times the Float64 CPU time -- because it runs at a lower clock ($1.4$--$2.6$
against about $4$\,GHz) and issues its instructions in order, waiting for the
result of each, without the out-of-order execution of a CPU core. The first
threads then gain less than the CPU threads do: $16$ GPU threads are only
$7.4$--$7.5$ times faster than one (L4, A100, RTX PRO 6000). The threads of a
warp of $32$ share one instruction stream, and whenever one of them finishes a
point, the store of its result and the fetch of the next point are executed
while the others wait. Once the first warp is full, however, the time falls in
exact proportion to $1/n$ over the next three decades, from $64$ to some
$10^{4}$--$10^{5}$ threads (the dotted line through the RTX PRO 6000 in
Fig.~\ref{fig:parallel}), at a parallel efficiency
$t_1/(n\,t(n))=\ParGpuEtaLo$--$\ParGpuEtaHi$ relative to one thread. In
Float64 the curves flatten between $64$ and $256$ threads, where the single
thread block occupies one multiprocessor whose few double-precision units are
already saturated by two warps; the slope resumes as soon as the blocks spread
over the device. The curves level off once every multiprocessor holds as many
warps as its registers allow -- on the RTX PRO 6000 near its capacity of
$2.9\cdot10^{5}$ resident threads.

\subsection{A performance model}\label{sec:parallel-model}

The measurements condense into a model with two hardware numbers and one
efficiency. If one thread needs $t_1$ per point and the device sustains
$P_{\mathrm{eff}}=t_1/t_{\mathrm{sat}}$ concurrent marches, a chart of $N$
points on $n$ threads takes
\begin{equation}\label{eq:parmodel}
  T(N,n) \approx \frac{N\,t_1}{\min\{\eta\,n,\,P_{\mathrm{eff}}\}} ,
\end{equation}
with $\eta\approx1$ on a CPU up to its physical core count and
$\eta\approx\ParGpuEtaLo$--$\ParGpuEtaHi$ on the GPUs in Float32 for $n\ge64$.
On a CPU, $P_{\mathrm{eff}}$ is the number of cores plus a fraction for
hardware threads. On the four GPUs it is $1.2$ to $2.6$ times the number of
CUDA cores in Float32 (Table~\ref{tab:parallel}). It exceeds the core count
because a lone thread is limited by instruction latency, not by throughput:
it waits for the result of each instruction, whereas a saturated
multiprocessor issues an instruction in every cycle from one of its many
resident warps. (The ratio also depends on how cores are counted: the T4 and
the A100 have integer units that are not counted as CUDA cores.) The march is
bound by instruction issue, not by memory: a point reads a few parameters and
writes some twenty bytes, about $55$\,GB/s on the fastest card, a few percent
of its memory bandwidth, and the L4 is $2.7$ times faster than the T4 at almost
the same bandwidth. The saturated cost times the core count,
$t_{\mathrm{sat}}\times\text{cores}$, varies by less than a factor of $1.5$
between the four cards, so the core count and a per-point cost measured once
estimate the chart time on another NVIDIA card within that factor; hardware of
other vendors, reachable through \pkg{KernelAbstractions.jl}, was not tested.
$P_{\mathrm{eff}}$ also carries over between problems: the two-mode turning
model, with $4.4$ times the evaluations of the fourth-order benchmark, each
$1.4$ times as expensive, needs $100\,\mu$s per point on one thread of the
RTX PRO 6000 and $3.0$\,ns per point saturated, an effective parallelism of
$33\,500$ against $\ParBestPeff$. For an interactive loop a fixed per-frame
cost -- kernel launch and the transfer of the result -- adds to
Eq.~\eqref{eq:parmodel}.

Float64 follows the same model with a larger $t_1$, and Fig.~\ref{fig:parallel}
shows the consequence of the hardware design. The graphics-class chips of the
T4, the L4 and the RTX PRO 6000 execute double-precision operations at
$1/32$--$1/64$ of the single-precision rate, so their saturated Float64 time is
$23$--$42$ times the Float32 time (only part of the instruction stream is
double precision), whereas the A100, built for scientific computing, pays a
factor of three. On this benchmark the Float32 chart equals the Float64 chart
except at a single flagged point of two million (Table~\ref{tab:gpu}), so
Float64 serves as a reference on the GPU rather than as a production format.
Float32 is adequate as long as the phase error of the delay terms,
$\epsilon_{32}\,\wmax\tau_{\max}$ ($3\cdot10^{-3}$\,rad here), stays far below
the step tolerance of the march; flagged points are re-checked in Float64.

\subsection{Real-time charts at 8K resolution}\label{sec:parallel-realtime}

\begin{table}[tb]
  \centering
  \caption{Brute-force charts of the fourth-order benchmark with the discrete
  march, one GPU thread per point, on four NVIDIA GPUs (Google Colab): median
  wall time of the kernel call in milliseconds, including launch and
  synchronization (full HD, $1920\times1080\approx2.07\cdot10^{6}$ points, ten
  repetitions; 8K, $7680\times4320\approx3.3\cdot10^{7}$ points, Float32,
  $\wmax=10^{5}$, five repetitions). $\wmax=10^{5}$ is the production window;
  $\wmax=15$ suffices for this model and is required by Float16, whose range
  ends at $65\,504$; in the Float16 column $\Dfun$ is evaluated in Float16 on a
  rescaled frequency axis while the march runs in Float32. The throughput
  column, in million points per second, refers to Float32. The timings of the
  T4 (passively cooled, 70\,W) scatter by up to a factor of two between cases
  under its thermal clock control -- its 8K chart costs half as much per point
  as its full-HD chart; at 8K the median exceeds the minimum by $18$--$28\,\%$
  on the A100 and the L4 as well.}\label{tab:gpu}
  \maybeinput{tables/tab_gpu.tex}
\end{table}

\begin{table}[tb]
  \centering
  \caption{The discrete march with $\Dfun$ evaluated in narrow number formats
  ($400\times400$ chart of the fourth-order benchmark, $\wmax=15$, best of four
  frequency scalings $\lam=\Omega\mu$ per format; the march itself runs in
  Float32). \emph{Flagged} points carry a sub-resolution or residual flag and
  are re-checked in a wider format; \emph{failed} points did not reach $\wmax$.
  $^\ast$Emulated: every operation is rounded to the format.}\label{tab:precision}
  \maybeinput{tables/tab_precision.tex}
\end{table}

The horizontal lines of Fig.~\ref{fig:parallel} translate the per-point cost
into frame rates. A chart refreshed thirty times per second leaves
$33$\,ms per frame; for a full-HD chart ($2.07\cdot10^{6}$ points) this needs
$16$\,ns per point, and for an 8K chart ($3.3\cdot10^{7}$ points) $1$\,ns per
point. A single CPU core is more than two orders of magnitude above the first
line, and even the $48$-thread server is a factor of six above it: a full-HD
chart takes $0.2$\,s, five frames per second, and thirty frames per second
would require about one sixth of the pixels ($800\times430$). All four GPUs
are below the full-HD line. Only the fastest, the \GpuBest{}, is below the 8K
line as well: at $\ParBestNsPerPoint$\,ns per point on the work queue of
Fig.~\ref{fig:parallel} ($\ParBestFullHDms$\,ms for full HD, $\ParBestEightKms$\,ms
for 8K), and with the production schedule of one thread per point at
$\GpuBestFthirtytwo$\,ms for the full-HD chart and $\GpuEightKFthirtytwo$\,ms for
the 8K chart in Float32 (Table~\ref{tab:gpu}). The other cards compute the 8K
chart in $40$--$104$\,ms. For comparison at equal precision, the Float64 march
on the \GpuBest{} is $\ParBestVsCpuFsixtyfour$ times faster than one Float64
CPU thread and $\ParBestVsServerFsixtyfour$ times faster than the whole
$48$-thread server; in Float32 it is $\ParBestVsServer$ times faster than the
server.

On this card the chart can therefore be recomputed \emph{while a slider is
being moved}: the delay, a gain or a damping ratio is changed, and the complete
chart -- at 8K more than thirty million stability decisions, each with an
estimate of the root nearest the imaginary axis -- is recomputed within the
time of one video frame. What then limits the frame rate is the transfer, not
the march. A loop that changes the delay for every frame and copies two result
fields, $\Zraw$ and $\sest$, into host memory runs at $\GpuBestSlider$\,ms per
frame in full HD but at $\GpuEightKSlider$\,ms (thirteen frames per second) at
8K, where copying several hundred megabytes dominates. An interactive
application keeps the result on the device and transfers only the colored
display image, whose coloring and transfer took about $2$\,ms per frame at 8K
(\code{gpu/results/interactive\_G4.csv}; these steps do not depend on the
march kernel).

These frame rates belong to the given problem. The cost per point is
proportional to the cost of one evaluation of $\Dfun$ and to the number of
evaluations, so a model whose $\Dfun$ is a large determinant or whose phase
march needs more steps moves all curves of
Fig.~\ref{fig:parallel} upwards by the corresponding factor: on the same card
the DAE showcase ($1.2$\,ns per point) runs at about $24$ frames per second at
8K, and the two-mode turning model ($3.0$\,ns per point) at about ten frames
per second at 8K and above $150$ at full HD. What does not change is the slope:
whatever the characteristic function, the chart time is proportional to $1/n$,
from a single core to tens of thousands of GPU threads.

\paragraph{Narrow number formats.} Since the march needs only the
\emph{phase} of $\Dfun$ to a fraction of a radian, the characteristic function
can even be evaluated in half precision (Table~\ref{tab:precision}). Float16
is accurate wherever no root is close to the line: after rescaling the
frequency into its range it yields, on the full-HD chart of
Table~\ref{tab:gpu}, $\GpuFsixteenWrong\,\%$ wrong counts that are \emph{all}
flagged, and $\GpuFsixteenFlagged\,\%$ flagged points in total; re-checking the
flagged points in Float32 -- in place on the device: gather, march, scatter --
restores the exact chart. It does not, however, make the chart faster: at 8K
the Float16 evaluation takes $1.01$--$1.12$ times the Float32 time on all four
cards (the factor of three of the T4 at full HD is an artifact of its clock
control). Scalar half-precision instructions issue at the single-precision
rate -- a gain would require packing pairs of values into vector registers --
the elementary functions are evaluated by promotion to Float32, and the
conversions have a reduced throughput of their own. (An earlier version of the
kernel, which passed the dual numbers through local memory, did run three
times faster in Float16, \code{gpu/results/interactive\_G4.csv}; the gain came
from the smaller memory traffic and disappeared once that traffic was
removed.) Float32 is therefore the production format on every card measured,
and the value of the narrow formats lies in what they reveal about the method:
how little precision the phase needs. BFloat16, with only seven mantissa bits,
flags $14\,\%$ of the chart, and the 8-bit and 4-bit formats are useless for
this purpose: with one to three mantissa bits even the \emph{parameters} of a
chart are quantized to a handful of distinct values per octave, so that
$95$--$100\,\%$ of the points are flagged and up to one in ten counts is
wrong; Float4 gets every count wrong. Load balancing, finally, does not pay:
the dynamic work queue of Fig.~\ref{fig:parallel}, in which threads that finish
early fetch new points, is $1.3$--$1.9$ times \emph{slower} than one thread per
point on every card, because the cost per point varies only moderately and the
hardware scheduler absorbs the difference at a lower overhead than the queue.
